\subsection{CONTINUOUS REAL-VALUED FUNCTIONS}

Besides the general properties of continuous functions with values in an arbitrary topological space (Chapter I, § 2), continuous real-valued functions have the following two fundamental properties:

THEOREM 1 (Weierstrass). Let \(f\) be a continuous real-valued function defined on a non-empty quasi-compact space \(\mathbf{X}\) . Then there is at least one point \(a \in \mathbf{X}\) such that \(f(a) = \sup_{x \in \mathbf{X}} (f(x))\) , and at least one point \(b \in \mathbf{X}\) such that \(f(b) = \inf_{x \in \mathbf{X}} f(x)\) . For \(f(\mathbf{X})\) is compact (Chapter I, § 9, no. 4, Theorem 2) and therefore closed in \(\overline{\mathbf{R}}\) ; hence \(f(\mathbf{X})\) contains its bounds.

This theorem is often stated in the form that a continuous real-valued function on a non-empty quasi-compact space attains its bounds.

COROLLARY. If a real-valued function defined on a non-empty quasi-compact space X is continuous and finite on X, then it is bounded in X.

THEOREM 2 (Bolzano). Let \(f\) be a continuous real-valued function defined on a connected space \(\mathbf{X}\) . If \(a\) and \(b\) are any two points of \(\mathbf{X}\) , and if \(\alpha\) is a real number belonging to the closed interval whose end-points are \(f(a)\) and \(f(b)\) , then there is at least one point \(x \in \mathbf{X}\) such that \(f(x) = \alpha\) .

For \(f(\mathbf{X})\) is connected (Chapter I, § 11, no. 2, Proposition 4) and is therefore an interval of \(\overline{\mathbf{R}}\) ( \(\S 4\) , no. 2, Proposition 5); hence it contains the closed interval with end-points \(f(a)\) and \(f(b)\) .

This theorem is often expressed in the form that a continuous real-valued function on a connected space cannot pass from one value to another without passing through all intermediate values.

This property is by no means characteristic of continuous functions; there are examples of functions defined on a connected space and discontinuous at every point which have this property (Exercise 2).

